\section{その他}

\begin{shortanswer}[年金1 2021 問題2(3)]
  公的年金に関し、「年金制度の機能強化のための国民年金法等の一部を改正する法律（令和2 年
法律第40 号） 」による年金受給開始時期の選択肢の拡大に関連するものとして、次の(1)および(2)
の設問に答えなさい。
\begin{enumerate}[(1)]
  \item 2022 年（令和4 年）4 月より施行されることとなった、老齢厚生年金および老齢基礎年金の
受給開始時期の選択肢に関する見直しの内容について簡記しなさい。
  \item 2023 年（令和5 年）4 月より施行されることとなった、70 歳到達後に老齢厚生年金および老
齢基礎年金の請求を行い、かつ請求時点で繰下げ受給を選択しない場合の取り扱いに関する見
直しの内容について簡記しなさい。
\end{enumerate}

\end{shortanswer}

\begin{answerbox}
  \begin{enumerate}[(1)]
  \item 
  \begin{itemize}
    \item 繰下げ受給の上限年齢が70歳から75歳へ見直される。繰下げ増額率は70歳までと変わらず、1月あたり0.7\%で、最大75歳まで繰り下げれば最大84\%となる。
    \item 繰上げ受給については平均余命の延伸に伴い、1月あたりの減額率が0.5\%から0.4\%となる。
  \end{itemize}

  \;

  \item 
  \begin{itemize}
    \item 70歳到達後に年金の請求を行う場合において、繰下げ受給を選択せずに、本来の支給開始年齢からの年金を選択した場合は、請求時から遡って5年を超える期間分の年金は時効により消滅していた。
    \item 改正後は、70歳到達後80歳到達前に年金の請求を行う場合において、繰下げ受給を選択せずに、遡って通常受給を選択した場合は、請求の5年前に繰下げ受給の選択があったものとみなし、受給権発生から請求の5年前までの月数に応じて増額された年金額が支給される
  \end{itemize}
\end{enumerate}
\end{answerbox}

\begin{explanationbox}
  過去の時事問題
\end{explanationbox}

\begin{shortanswer}[年金1 2022 問題2(3)]
  公的年金に関し、次の(1)、(2)の設問に解答しなさい。 

  \begin{enumerate}[(1)]
    \item 新型コロナウイルス感染症の影響に伴う休業により著しく報酬が下がった場合における、標
準報酬月額の特例改定のうち、2022（令和4 年）年8 月または2022 年（令和4 年）9 月に
新たに休業により著しく報酬が下がった者の特例措置の内容および対象者の要件について簡
潔に入力しなさい。
  \item 2022 年（令和4 年）4 月より実施された加給年金の見直しの内容について簡潔に入力しなさ
い。ただし、経過措置の内容については触れなくて良いものとする。
  \end{enumerate}

\end{shortanswer}

\begin{answerbox}
    \begin{enumerate}[(1)]
    \item ＜特例措置の内容＞
    
事業主から届出があった場合に、休業により報酬が著しく低下した月に受けた報酬の総額を報
酬月額として算定し、翌月から標準報酬月額を改定できる取扱いとするもの。

＜対象者の要件＞

以下の要件をすべて満たしていること。
\begin{itemize}
  \item 新型コロナウイルス感染症の影響による休業があったことにより、令和 4 年 8 月または
令和4 年9 月の報酬が著しく下がった者
\item 著しく報酬が下がった月に支払われた報酬の総額（1 か月分）が、既に設定されている標
準報酬月額に比べて2 等級以上下がった者
\item 本特例措置による改定内容に書面により同意している者
\end{itemize}
  \;

  \item 加給年金の加算対象となる配偶者が、被保険者期間が20年（中高齢者等の特例に該当する方を含む）以上ある老齢、退職を支給事由とする年金の受給権を有する場合、配偶者に対する給付が全額支給停止されているときには、加給年金が支給されることとなっていたが、見直しにより加給年金が支給停止されることとなった。
  \end{enumerate}

\end{answerbox}

\begin{explanationbox}
  過去の時事問題
\end{explanationbox}

\begin{fillblank}[年金1 H28 問題1(2)J]
  平成 24 年に可決・成立した年金機能強化法では、年金の受給資格期間の（　　A　　）年への短縮が
規定されている。
\end{fillblank}

\begin{answerbox}
  \begin{enumerate}[A: ]
    \item 10
  \end{enumerate}
\end{answerbox}

\begin{shortanswer}[数理人試験 2024 問題3 設問3]
  令和4年10月1日以降に開始した育児休業等の期間について、社会保険料が免除となる要件を簡記せよ。
\end{shortanswer}

\begin{answerbox}
  毎月の給与に係る社会保険料については、育休開始日が属する月から、終了日の翌月が属する月の前月までの保険料に加え、育休開始日の属する月内に14日以上の育児休業を取得した場合も当該つきの社会保険料が免除される。

また、賞与に係る社会保険料については、賞与を受け取った月の末日を含む、連続した1ヶ月を超える育児休業を取得した場合に限り社会保険料が免除される
\end{answerbox}

\begin{shortanswer}
  短期間労働者への公的年金の適用においては、所謂、「年収の壁」に収まるように就業を調整す
ることに対する対応策が議論されている。以下の2つの「年収の壁」についてそれぞれ概要及び
判定基準の算定対象となる給与の種類を簡記せよ。
\begin{enumerate}[(1)]
  \item 年収130万円の壁
  \item 年収106万円の壁
\end{enumerate}
\end{shortanswer}

\begin{answerbox}
  \begin{enumerate}[(1)]
  \item 年収130万円を超えると社会保険の被扶養者（第三号被保険者）から外れることとなる判定基準の算定対象となる給与：前年度の所得証明書に基づき賞与等を含めた全収入ベースで判定される
  \item 月8.8万円（＝年収106万円）未満の場合、被用者年金の対象要件の適用外となる判定基準の算定対象となる給与：基本給及び諸手当（地域手当等）を基に判定し、それ以外の通勤手当、賞与等は含まない
\end{enumerate}
\end{answerbox}

\begin{explanationbox}
「年収の壁」への対応

  \url{https://www.mhlw.go.jp/stf/taiou_001_00002.html}
\end{explanationbox}

\begin{fillblank}[年金1 2025 問題1(5)ア]
  年金積立金は、 （　　A　　） の寄託を受けた （　　B　　） が運用している。2025 年 4 月から適用
されている基本ポートフォリオにおける資産構成割合は国内株式が （　　C　　） ％、外国株式が
（　　D　　） ％である。
\end{fillblank}

\begin{answerbox}
  \begin{enumerate}[A: ]
  \item 厚生労働大臣
  \item 年金積立金管理運用独立行政法人（GPIF）
  \item 25
  \item 25
\end{enumerate}
\end{answerbox}